% Inactive reconstruction of the continuous-solution regime, 2 <= m < n.
% Deep inputs are declared below. General bounded-solution continuity is separate.
\section{Smoothness of continuous admissible solutions}

Let $(X,\omega)$ be compact Hermitian, of dimension $n$, and let
$2\leq m<n$. A continuous real function $v$ is admissible if
\[
 (\omega+\ddc v)\wedge\gamma_1\wedge\cdots\wedge\gamma_{m-1}
                              \wedge\omega^{n-m}\geq0
\]
in distributions for every smooth collection of forms in
$\Gamma_m(\omega)$. For smooth $v$ this is equivalent to the pointwise
closed-cone condition. We write
$H_m(v)=(\omega+\ddc v)^m\wedge\omega^{n-m}$.

The argument uses two established continuous pluripotential inputs:
the continuous Hessian products constructed by uniform smooth admissible
approximation, with continuity in that topology (Ko\l odziej--Nguyen,
1507.06755v2, Proposition 2.9 and Corollary 2.10), and continuous strict
domination (Guedj--Lu, 2302.03354v1, Corollary 2.3). The latter says that
\[
 \mathbf1_{\{u<v\}}H_m(u)\leq c\mathbf1_{\{u<v\}}H_m(v),\quad 0\leq c<1,
 \qquad\Longrightarrow\qquad u\geq v.
\]
In particular equality of $e^{-u}H_m(u)$ and $e^{-v}H_m(v)$ implies
$u=v$: apply strict domination on $\{u<v-\delta\}$ with
$c=e^{-\delta}$, then let $\delta\downarrow0$, and exchange the functions.
These are inputs for continuous functions, not newly recertified general
bounded-product or bounded-domination theorems. Their selected primary
proof scopes and deeper algebraic/envelope dependencies remain explicit
in the audit.

For higher regularity we use the classical Evans--Krylov--Caffarelli
estimate in Tosatti--Wang--Weinkove--Yang, 1402.0554v2, Theorem 3.4,
including its continuous-viscosity version stated immediately after the
theorem. The cone estimates, smooth trace estimate and gradient blowup
are the preceding reconstructed fragments 599, 600 and 596. Smooth
elliptic inversion, the distributional linear maximum principle and
Schauder bootstrapping are standard analytic inputs.

\begin{lemma}\label{draft:604-smooth-proper-existence}
For every $a>0$ and smooth real $H$ on $X$, there is a unique smooth
strictly admissible solution of
\[
 H_m(w)=e^{aw+H}\omega^n.
\]
It satisfies
\[
 -\frac{\max_X H}{a}\leq w\leq-\frac{\min_X H}{a}.
\]
\end{lemma}
\begin{proof}
At a maximum of $w$, the admissible form $\omega+\ddc w$ is at most
$\omega$, hence $e^{aw+H}\leq1$ there. At a minimum it is at least
$\omega$, and the reverse inequality holds. These give the bounds.
For two solutions, a positive maximum of their difference would imply
$H_m(w_1)\leq H_m(w_2)$ there, contradicting the exponential factor.
This proves uniqueness.

First take $a=1$ and use the path
$H_m(w_t)=e^{w_t+tH}\omega^n$, starting with $w_0=0$.
The trace theorem 600 applies with prescribed potential $v=0$,
fixed smooth comparison $0$, and $h(x,v)=tH(x)$; its constants are
uniform for $t\in[0,1]$. The gradient proposition 596 then gives a
uniform gradient and complex-Hessian bound. The interior-Laplace argument
of 603 gives uniform $W^{2,p}$ and hence $C^{1,\alpha}$ bounds, for any
fixed $\alpha<1$. Consequently
$r_t=e^{(w_t+tH)/m}$ has a uniform H\"older bound.

The extension of 603 is uniformly elliptic and concave on all real
Hessians, with coefficient oscillation bounded by
$C|x-y|(1+\|R\|)$. It therefore belongs to the class in Theorem 3.4
of the stated analytic input (trace and operator-norm ellipticity differ
only by the fixed factor $2n$). That theorem gives uniform $C^{2,\beta}$
bounds on a finite cover. The original smooth Hessian operator is
uniformly elliptic near the attained compact cone set. Differentiating
and applying Schauder estimates gives bounds for every higher derivative.

The linearization of
$\log(H_m(w)/\omega^n)-w-tH$ is $L_w-1$, with positive definite
principal coefficients. The maximum principle gives zero kernel;
smooth elliptic Fredholm theory gives surjectivity and Schauder inversion.
The implicit-function theorem proves openness. The uniform estimates
prove closedness, so the path reaches $t=1$.
For general $a$, use the metric $a\omega$ and potential $\widehat w=aw$.
Both sides of its coefficient-one equation acquire the same factor
$a^n$, so it is exactly the stated equation after division by that factor.
\end{proof}

\begin{lemma}\label{draft:604-continuous-smooth-approximation}
Every continuous admissible $\varphi$ on $X$ is the uniform limit of
smooth strictly admissible functions.
\end{lemma}
\begin{proof}
Fix a smooth $h\geq\varphi$ and a constant
\[
 b>\max\left\{1,\left\|H_m(h)/\omega^n\right\|_\infty\right\}.
\]
For each positive integer $N$ solve, by the preceding lemma,
\[
 H_m(w_N)=b e^{N(w_N-h)}\omega^n.
\]
At a positive maximum of $w_N-h$, the form $\omega+\ddc h$ lies above
the strictly admissible form $\omega+\ddc w_N$, hence is admissible there.
Monotonicity would give
$b e^{N(w_N-h)}\leq H_m(h)/\omega^n<b$, a contradiction.
Thus $w_N\leq h$. The preceding zero-order estimate also gives
\[
 \|w_N\|_\infty\leq M_N:=\|h\|_\infty+\frac{\log b}{N}.
\]

Let $x_N$ maximize $\varphi-(1+N^{-1})w_N$.
The smooth positive form
$T_N=(\omega+\ddc w_N)^{m-1}\wedge\omega^{n-m}$
defines a strictly elliptic linear operator. Admissibility gives
$(\omega+\ddc\varphi)\wedge T_N\geq0$ in distributions.
Its linear maximum principle, applied to the upper test
$(1+N^{-1})w_N$ plus the maximum value, gives at $x_N$
\[
 0\leq\bigl(\omega+(1+N^{-1})\ddc w_N\bigr)\wedge T_N.
\]
This application is to a continuous distributional subsolution of a
smooth strictly elliptic operator. It does not assume a second derivative
of $\varphi$. Writing $d_N=H_m(w_N)/\omega^n$, polarized G\aa rding
positivity yields
\[
 (N+1)d_N\geq \omega\wedge T_N/\omega^n
                         \geq d_N^{(m-1)/m}.
\]
Therefore, at this maximum point,
\[
 w_N-h\geq-E_N,\qquad
 E_N=\frac{m\log(N+1)+\log b}{N}.
\]
Since $\varphi\leq h$, the maximum of
$\varphi-(1+N^{-1})w_N$ is at most $E_N+M_N/N$.
It follows on all of $X$ that
\[
 \varphi-E_N-2M_N/N\leq w_N\leq h.
\]
Given $\delta>0$, choose an ordinary smooth majorant
$\varphi\leq h\leq\varphi+\delta$, and then $N$ so large that
$E_N+2M_N/N\leq\delta$. The resulting $w_N$ is strictly admissible
and within $\delta$ of $\varphi$. Let $\delta\downarrow0$.
No monotonicity of the chosen approximants is needed.
\end{proof}

\begin{proposition}\label{draft:604-proper-uniform-stability}
Let $H_j$ be smooth and converge uniformly to a continuous $H$.
The smooth solutions $w_j$ of
$H_m(w_j)=e^{w_j+H_j}\omega^n$ converge uniformly to the unique
continuous admissible solution of
$H_m(w)=e^{w+H}\omega^n$. For smooth $H_1,H_2$,
\[
 \|w_1-w_2\|_\infty\leq\|H_1-H_2\|_\infty.
\]
\end{proposition}
\begin{proof}
At a maximum of $w_1-w_2$, monotonicity gives
$e^{w_1+H_1}\leq e^{w_2+H_2}$, hence the maximum is at most
$\|H_1-H_2\|_\infty$. Exchange the functions to obtain the norm bound.
Thus $w_j$ is Cauchy uniformly. The limit is admissible by passage of
the linear distributional inequalities to the limit. The continuous
Hessian-product input passes the equation to this limit. The exponential
order consequence stated above proves uniqueness among continuous
admissible solutions, independently of the approximation of $H$.
\end{proof}

\begin{theorem}\label{draft:604-continuous-hessian-regularity}
If $\varphi$ is continuous and admissible on $X$, $f(x,t)$ is smooth
on a neighborhood of its bounded range, and
\[
 H_m(\varphi)=e^{f(x,\varphi)}\omega^n,
\]
then $\varphi$ is smooth.
\end{theorem}
\begin{proof}
Choose smooth admissible $v_j\to\varphi$ uniformly, by the approximation
lemma. Their values lie in a fixed compact interval on which $f$ is
smooth. Put $h(x,t)=-t+f(x,t)$ and solve
\[
 H_m(u_j)=e^{u_j+h(x,v_j)}\omega^n.
\]
The preceding stability proposition shows that $u_j$ converges uniformly
to a continuous admissible $u$ satisfying
\[
 e^{-u}H_m(u)=e^{-\varphi+f(x,\varphi)}\omega^n
                              =e^{-\varphi}H_m(\varphi).
\]
Continuous exponential uniqueness gives $u=\varphi$.

The zero-order bounds for $u_j$ are uniform. Fix a smooth ordinary
approximation $\widetilde\varphi$ within the small tolerance required
by the trace theorem 600. After discarding finitely many terms, all
$v_j$ lie within that fixed tolerance of $\widetilde\varphi$.
The trace theorem and gradient blowup 596 then give uniform gradient
and complex-Hessian bounds for $u_j$, independent of derivatives of
$v_j$. The first lemma of 603 yields
$\varphi\in W^{2,p}$ for every finite $p$, and
$\varphi\in C^{1,\alpha}$ for every $\alpha<1$.

The last proposition of 603 makes $\varphi$ a continuous viscosity
solution of the concave uniformly elliptic extension, with right-hand
side $e^{f(x,\varphi)/m}$. This right-hand side is H\"older. The
continuous-viscosity version of Theorem 3.4 therefore gives
$\varphi\in C^{2,\beta}$, for some $\beta>0$.
The operator agrees with the original Hessian root on the attained
compact cone set; the now classical equation thus has smooth dependence
on $x,\varphi$ and its complex Hessian, and is uniformly elliptic there.
Difference quotients and interior Schauder estimates give $C^{3,\beta}$;
iteration proves smoothness. A second family of auxiliary equations is
unnecessary because the analytic estimate applies to the viscosity limit.
\end{proof}

This theorem assumes continuity of the given weak solution. Extending
it to an arbitrary bounded admissible weak solution requires its bounded
Hessian products, domination and continuity to be justified separately.
No Orlicz statement or conclusion for a merely semipositive background
is asserted here.
